\scanpage{90}

``Invariants'' des groupes des polyèdres réguliers sphériques

\fbox{1} Soit $S$ schéma de base, $(\unsure{E|F},q,\omega)$ un Module spécial quadratique sur $S$ de rang $3$, donc
\[
\delta(q)=-\omega^{\unsure{\otimes(2)|(\otimes2)|\otimes2}}\qquad\omega\in\Gamma(S,(\underline{\det}E)^*)\tag*{(1)}
\]

On sait alors que $E$ est muni d'une loi $[\, ,\, ]$, plus précisément, on peut trouver une algèbre de quaternions tordus $\mathcal M$ sur $S$ (définie à isom unique près) et un isomorphisme de $\underline O_S$-Mod.
\[
E\xleftarrow{\sim}\unsure{\gamma|\mathscr V}\mathcal M=\operatorname{Ker}\bigl(\mathcal M\xrightarrow[\text{trace \uline{réduite}}]{\mathrm{tr}}\underline O_S\bigr)\tag*{(2)}
\]
transformant $q$ en la forme
\[
q_{\mathcal M}(u)=-\det u\qquad(u\in\Gamma(-,\gamma\mathcal M))\tag*{(\unsure{3|B})}
\]
et $\omega$ en la forme \inserted{$\omega_{\mathcal M}$} définie par la condition que
\[
\varphi_{\mathcal M}([u,v],w)\omega_{\mathcal M}=2u\wedge v\wedge w\qquad(u,v,w\in\Gamma(-,\gamma\mathcal M))\tag*{(4)}
\]
% slanted
\margin{on définit en tte rigueur $\varphi_{\mathcal M}$ que \unsure{si|qd} $2$ inv.}
où $\varphi_{\mathcal M}$ est la forme \deleted{\unsure{quadratique}} \inserted{bilinéaire} sym. associée à $q_{\mathcal M}$
\[
\left\{\begin{aligned}q_{\mathcal M}(u+v)&=\deleted{\det}q_{\mathcal M}(u)+q_{\mathcal M}(v)+\varphi_{\mathcal M}(u,v)\\ \text{i.e.}\quad\det(u+v)&=\det(u)+\det(v)-\varphi_{\mathcal M}(u,v)\end{aligned}\right.\tag*{(5)}
\]
pour $u,v\in\Gamma(-,\gamma\mathcal M)$, ce qui donne aussi
\[
\varphi_{\mathcal M}(u,v)=\operatorname{Tr}u\circ v\qquad u,v\in\Gamma(-,\gamma\mathcal M)\tag*{(6)}
\]
% archive sheet number: 89
